\section{Defenses}
\label{sec:defense}

\subsection{Why the obvious controls fail}
\label{sec:matrix}

Table~\ref{tab:defense} collects the controls a defender would reach for, including the
one this paper proposed and then falsified. The pattern is uncomfortable: of six
candidate controls, four fail outright, one is partial and cost-dependent, and the only
reliably effective intervention is removing the channel altogether---which also removes
whatever legitimate coordination it was carrying.

The reason is visible in P1, P4 and the throughput probe together. The channel effect
\textbf{saturates on almost nothing}: a two-item menu, a single $60$-character message,
or three of twelve rounds of access each recover most of the effect of an unrestricted
channel. Controls that ration capacity therefore operate on the wrong margin, and
controls that inspect content operate on a signal \S\ref{sec:horserace} shows is not
predictive.

\begin{table}[t]
\centering
\caption{\textbf{The defense matrix.} Established results are from prior work and
\S\ref{sec:reanalysis}; the last two rows are this paper's constructive contribution.}
\label{tab:defense}
\small
\begin{tabularx}{\textwidth}{l X l}
\toprule
control & result & verdict \\
\midrule
format restriction (menu / intention-only) & lock-in $\to 0.00$--$0.06$ & works, but
destroys legitimate coordination and arms deception \\
throughput throttling & one $60$-char SMS: $0.45\!\to\!0.80$ & \textbf{fails} \\
passive monitoring & detect $0.85$--$1.00$, $\Delta K\approx0$ & \textbf{fails} \\
auditing the report channel & $0.71$ unobs.\ vs $0.77$ obs. & \textbf{fails} \\
\midrule
calibrated monitor $+$ enforcement & loss $-64$ to $-85\%$; useful in $14/24$ regimes
& \textbf{partial, bounded} \\
opening-window suppression (P4) & $0.53\times$ linear; partial suppression $\approx$
full channel & \textbf{fails} \\
\bottomrule
\end{tabularx}
\end{table}

\subsection{A cost-optimal monitor operating point}
\label{sec:monitor}

The prior work reported the collusion monitor as discriminative but mis-thresholded.
That is a diagnostic, not a defense: it does not tell an operator where to cut. We
turn it into a decision rule. The principal penalises on a flag, so per scored round
the expected loss at threshold $t$ is
\[
L(t) \;=\; c_{\text{FP}}\,(1-\pi)\,\mathrm{FPR}(t) \;+\; c_{\text{FN}}\,\pi\,\mathrm{FNR}(t),
\]
with $\pi$ the prevalence of colluding rounds. The operator picks $t^\star =
\arg\min_t L(t)$ and---critically---compares $L(t^\star)$ against the two do-nothing
baselines, never-flag ($c_{\text{FN}}\pi$) and always-flag ($c_{\text{FP}}(1-\pi)$).

Reproducing the published labelling (realised collusion, ambiguous band dropped) gives
$\mathrm{AUC}=0.893$. Against the shipped operating point (TPR $0.913$, FPR $0.450$),
calibration cuts expected loss by $64$--$85\%$ at $\pi{=}0.25$. But the honest result
is the boundary: \textbf{the monitor beats both do-nothing baselines in only $14$ of
the $24$ cost/prevalence regimes examined.} At high false-positive cost and low
prevalence the optimal policy is not to flag at all, and the correct operator
recommendation is to run no detector rather than a mis-thresholded one.

\begin{figure}[t]
\centering
\includegraphics[width=0.95\textwidth]{figures/p3_monitor_cost.pdf}
\caption{\textbf{Operating the monitor as a decision, not a diagnostic.} Left: the
score is discriminative. Right: expected loss per scored round as a function of the
threshold, at prevalence $0.25$, for three false-positive/false-negative cost ratios.
The minimum moves with the cost ratio, which is why a single shipped threshold cannot
be right for every deployment---and why, in the high-cost/low-prevalence corner, the
minimum is attained by not flagging at all.}
\label{fig:monitor}
\end{figure}

\subsection{Opening-window suppression}
\label{sec:p4}

\S\ref{sec:decomp} showed the channel both displaces the opening and slows the decay.
If the opening is where the trajectory is set, an operator can close the channel for
the first few rounds only and capture most of the protection while leaving agents free
to communicate for the rest of the interaction. P4 apportions the two components with
four arms---\texttt{full}, \texttt{suppress\_early} (rounds $0$--$2$ closed),
\texttt{suppress\_late} (only rounds $0$--$2$ open), \texttt{none}---at $30$ seeds
each, and reports protection per round suppressed so the control can be priced.

\paragraph{Result: the defense does not work.}
Neither window carries the effect, because \emph{partial suppression barely suppresses}
(Table~\ref{tab:p4}). Closing the channel for the opening rounds leaves lock-in
statistically indistinguishable from a fully open channel (RD $-0.067$
$[-0.225,+0.082]$); closing everything \emph{but} the opening is likewise
indistinguishable ($-0.100$ $[-0.266,+0.056]$). Only removing the channel outright
moves the outcome (\texttt{full} $-$ \texttt{none} $=+0.500$ $[+0.282,+0.667]$,
$p{=}2\times10^{-5}$).

\begin{table}[t]
\centering
\caption{\textbf{P4, $n{=}30$/arm.} Opening window $=$ rounds $0$--$2$. Every arm with
\emph{any} channel access lands near the full-channel ceiling; only total removal
separates.}
\label{tab:p4}
\small
\begin{tabular}{lccccc}
\toprule
arm & lock-in & Wilson 95\% CI & mean $C$ & opening $C$ & closing $C$ \\
\midrule
none            & 0.467 & [0.30, 0.64] & 0.718 & 0.911 & 0.577 \\
suppress\_early & 0.900 & [0.74, 0.97] & 0.936 & 0.964 & 0.931 \\
suppress\_late  & 0.867 & [0.70, 0.95] & 0.916 & 1.000 & 0.860 \\
full            & 0.967 & [0.83, 0.99] & 0.972 & 1.000 & 0.957 \\
\bottomrule
\end{tabular}
\end{table}

Priced as a control, suppressing three of a median twelve rounds costs $25\%$ of the
interaction and buys $13\%$ of the available protection---$\mathbf{0.53\times}$ the
naive linear rate, i.e.\ \emph{worse} than closing the channel uniformly. We proposed
this defense in the design and report it as \textbf{failed}.

\paragraph{Why it failed, and what that implies.}
The observational decomposition (\S\ref{sec:decomp}) was not wrong---the channel really
does both lift the opening and slow the decay---but it does not license the inference
we drew from it. The effect saturates on very little channel access: a pair that can
talk for nine of twelve rounds behaves like a pair that can talk throughout, and a pair
that can talk for only the first three behaves the same way. This converges with the
throughput probe, where a single $60$-character message moved a pair most of the way
(\S\ref{sec:sms}), and with P1, where a two-item menu was sufficient. \textbf{Rationing
a channel is not a control at any granularity we tested}; the operative choice is
whether the channel exists at all.

% ======================================================================
